\documentclass[10pt,conference]{IEEEtran}
\usepackage[utf8]{inputenc}
\usepackage{graphicx}
\usepackage{booktabs}
\usepackage{url}
\usepackage{balance}
\usepackage{xcolor}
\usepackage{multicol}

\begin{document}

%\title{Easier AI: If Anyone Can Do It for a Dollar,\\ Can You Do It for a Cent?}
\title{A Sweeter Lesson?\\ A Second Look at Big AI for Software Engineering}

% Double-anonymous submission: author block and acknowledgements are
% commented out until camera-ready.
\author{\IEEEauthorblockN{Anonymous Author(s)}\IEEEauthorblockA{Anonymous Institution}}
%\author{\IEEEauthorblockN{Tim Menzies}
%\IEEEauthorblockA{North Carolina State University, Raleigh, NC, USA \\ timm@ieee.org}}

\maketitle

\begin{abstract}
Sutton's bitter lesson says that general methods riding ever-cheaper
computation ultimately beat approaches built from human knowledge.
Written in 2019, before GPT-3, the essay predicted the scaling era we
now live in. This paper rereads it from inside software engineering
(SE) and reports a sweeter lesson. All four of Sutton's motivating
examples (chess, Go, speech, vision) enjoy fast, cheap, reliable
feedback, so computation can substitute for knowledge. Much of SE
analytics has no such oracle: labels cost hours of expert time or
whole test-suite runs. There, the bitter lesson's preconditions fail,
and radical simplicity wins. Across 127 multi-objective SE
optimization tasks, an active learner that labels just 50 examples,
then checks five, reaches 85\% of the best known result, via trees of
four to nine attributes, in two hours on one laptop. Yet in a survey
of 229 papers on LLMs in SE, only 5\% compared against anything
simpler. We propose an agenda to fix that: baseline the simple; ask,
at review time, ``compared to what simpler?''; and teach easier AI
first, scaling up only when the simple fails.
\end{abstract}

\begin{IEEEkeywords}
active learning, search-based software engineering, green AI, empirical standards
\end{IEEEkeywords}

\section{Introduction}
In March 2019, before GPT-3 existed, Sutton published a short essay
called {\em The Bitter Lesson}~\cite{sutton19}. Figure~\ref{fig:sutton}
reprints it in full. Its claim: general methods that leverage
computation beat methods built from human knowledge; always,
eventually, and by a large margin. The essay was prescient. It
described the scaling era before that era arrived, and its
unlimited-CPU perspective now looms over most current thinking about
AI, including AI for software engineering (SE). This paper is a
second look at that lesson, from inside SE. We will argue that, for
much of this field, there is a sweeter lesson going unlearned.

The cost of not looking compounds. Lehman's second law warns that a
system grows ever more complex unless work is done to reduce
it~\cite{lehman80}. AI-for-SE is now such a system: models, training
sets, energy budgets, and invoices all grow by orders of magnitude,
and little work pushes the other way.

This paper pushes the other way. We ask two questions. First: how
much of AI-for-SE could be radically simpler? Recent evidence,
summarized in \S\ref{sec:how}, says much of it~\cite{ganguly26}.
Second: if simple works so well, why do so few papers check? In a
survey of 229 papers on LLMs in SE, only 5\% baselined against
anything simpler~\cite{menzies26}. We call that a methodological
error. We also think it is a fixable one.

One caveat, stated early: generative tasks (translation, code
summarization, natural-language interfaces) genuinely need large
models. But many SE tasks (configuration, tuning, defect prediction,
testing, cost estimation) have far simpler solutions, and this paper
is about those.

The experiments cited here appeared, or will appear, in prior papers
by Ganguly, Rayegan, Senthilkumar, Srinivasan, Menzies, and
colleagues~\cite{ganguly26, rayegan26, senthilkumar26, srinivasan26,
menzies26}. What is new here is the rereading of Sutton, the
synthesis, and the agenda: \S\ref{sec:why} argues the why,
\S\ref{sec:how} the how, \S\ref{sec:whynot} diagnoses why not, and
\S\ref{sec:future} plans the full-length empirical result.

\begin{figure*}[!t]
\centering
\begin{minipage}{0.94\textwidth}
\fontsize{6.6}{7.0}\selectfont\ttfamily
\setlength{\parskip}{2.5pt}\setlength{\parindent}{0pt}
\begin{multicols}{2}
\centerline{THE BITTER LESSON}
\centerline{Rich Sutton, March 13, 2019}
\medskip

The biggest lesson that can be read from 70 years of AI research is
that general methods that leverage computation are ultimately the most
effective, and by a large margin. The ultimate reason for this is
Moore's law, or rather its generalization of continued exponentially
falling cost per unit of computation. Most AI research has been
conducted as if the computation available to the agent were constant
(in which case leveraging human knowledge would be one of the only
ways to improve performance) but, over a slightly longer time than a
typical research project, massively more computation inevitably
becomes available. Seeking an improvement that makes a difference in
the shorter term, researchers seek to leverage their human knowledge
of the domain, but the only thing that matters in the long run is the
leveraging of computation. These two need not run counter to each
other, but in practice they tend to. Time spent on one is time not
spent on the other. There are psychological commitments to investment
in one approach or the other. And the human-knowledge approach tends
to complicate methods in ways that make them less suited to taking
advantage of general methods leveraging computation. There were many
examples of AI researchers' belated learning of this bitter lesson,
and it is instructive to review some of the most prominent.

{\color{blue!55!black}
In computer chess, the methods that defeated the world champion,
Kasparov, in 1997, were based on massive, deep search. At the time,
this was looked upon with dismay by the majority of computer-chess
researchers who had pursued methods that leveraged human understanding
of the special structure of chess. When a simpler, search-based
approach with special hardware and software proved vastly more
effective, these human-knowledge-based chess researchers were not good
losers. They said that ``brute force'' search may have won this time,
but it was not a general strategy, and anyway it was not how people
played chess. These researchers wanted methods based on human input to
win and were disappointed when they did not.

A similar pattern of research progress was seen in computer Go, only
delayed by a further 20 years. Enormous initial efforts went into
avoiding search by taking advantage of human knowledge, or of the
special features of the game, but all those efforts proved irrelevant,
or worse, once search was applied effectively at scale. Also important
was the use of learning by self play to learn a value function (as it
was in many other games and even in chess, although learning did not
play a big role in the 1997 program that first beat a world champion).
Learning by self play, and learning in general, is like search in that
it enables massive computation to be brought to bear. Search and
learning are the two most important classes of techniques for
utilizing massive amounts of computation in AI research. In computer
Go, as in computer chess, researchers' initial effort was directed
towards utilizing human understanding (so that less search was needed)
and only much later was much greater success had by embracing search
and learning.

In speech recognition, there was an early competition, sponsored by
DARPA, in the 1970s. Entrants included a host of special methods that
took advantage of human knowledge---knowledge of words, of phonemes,
of the human vocal tract, etc. On the other side were newer methods
that were more statistical in nature and did much more computation,
based on hidden Markov models (HMMs). Again, the statistical methods
won out over the human-knowledge-based methods. This led to a major
change in all of natural language processing, gradually over decades,
where statistics and computation came to dominate the field. The
recent rise of deep learning in speech recognition is the most recent
step in this consistent direction. Deep learning methods rely even
less on human knowledge, and use even more computation, together with
learning on huge training sets, to produce dramatically better speech
recognition systems. As in the games, researchers always tried to make
systems that worked the way the researchers thought their own minds
worked---they tried to put that knowledge in their systems---but it
proved ultimately counterproductive, and a colossal waste of
researcher's time, when, through Moore's law, massive computation
became available and a means was found to put it to good use.

In computer vision, there has been a similar pattern. Early methods
conceived of vision as searching for edges, or generalized cylinders,
or in terms of SIFT features. But today all this is discarded. Modern
deep-learning neural networks use only the notions of convolution and
certain kinds of invariances, and perform much better.
}

This is a big lesson. As a field, we still have not thoroughly learned
it, as we are continuing to make the same kind of mistakes. To see
this, and to effectively resist it, we have to understand the appeal
of these mistakes. We have to learn the bitter lesson that building in
how we think we think does not work in the long run. The bitter lesson
is based on the historical observations that 1) AI researchers have
often tried to build knowledge into their agents, 2) this always helps
in the short term, and is personally satisfying to the researcher, but
3) in the long run it plateaus and even inhibits further progress, and
4) breakthrough progress eventually arrives by an opposing approach
based on scaling computation by search and learning. The eventual
success is tinged with bitterness, and often incompletely digested,
because it is success over a favored, human-centric approach.

One thing that should be learned from the bitter lesson is the great
power of general purpose methods, of methods that continue to scale
with increased computation even as the available computation becomes
very great. The two methods that seem to scale arbitrarily in this way
are search and learning.

The second general point to be learned from the bitter lesson is that
the actual contents of minds are tremendously, irredeemably complex;
we should stop trying to find simple ways to think about the contents
of minds, such as simple ways to think about space, objects, multiple
agents, or symmetries. All these are part of the arbitrary,
intrinsically-complex, outside world. They are not what should be
built in, as their complexity is endless; instead we should build in
only the meta-methods that can find and capture this arbitrary
complexity. Essential to these methods is that they can find good
approximations, but the search for them should be by our methods, not
by us. We want AI agents that can discover like we can, not which
contain what we have discovered. Building in our discoveries only
makes it harder to see how the discovering process can be done.
\end{multicols}
\end{minipage}
\caption{Sutton's bitter lesson, reprinted verbatim~\cite{sutton19}.
Blue: the essay's four motivating examples. Each comes from a domain
with a fast, cheap, reliable feedback oracle. \S\ref{sec:how} argues
that much of SE analytics has no such oracle.}
\label{fig:sutton}
\end{figure*}

\section{Why Easier AI?}
\label{sec:why}
Before the economics, reread the blue paragraphs of
Figure~\ref{fig:sutton}. Sutton offered four examples: chess, Go,
speech recognition, and computer vision. Note what they share. A
chess or Go program scores a candidate move by playing millions of
games against itself; speech and vision systems train against
millions of labeled utterances and images. In every case, feedback is
immediate, cheap, and incorruptible, so computation can substitute
for knowledge without limit. That is the essay's silent premise. We
do not dispute Sutton's history. We dispute the transfer: where
feedback is slow, dear, or scarce (\S\ref{sec:how}), scale loses its
lever, and the lesson turns sweeter.
Start with money. Thirty years ago, Wirth complained that software
gets slower faster than hardware gets faster~\cite{wirth95}.
Even as the price per token falls,
Jevons' paradox applies: cheaper inference means more inference, and
total token consumption has grown by orders of magnitude since 2023.
Bain \& Company forecasts that by 2030, the AI industry must find
\$800 billion per year in new revenue to fund its own compute
build-out; capital expenditure was near \$725 billion in 2026 and is
projected to pass \$1 trillion in 2027~\cite{bain25}. Someone must
pay for all this, and researchers who build on rented frontier models
should ask what happens when the subsidies stop.

Money is one worry. Trust is another. A model small enough to read is
a model that can be audited, repaired, and defended in front of a
regulator. Rudin argues that high-stakes decisions need interpretable
models, not post-hoc explanations of black boxes~\cite{rudin19}.
Bender et al.\ ask, of language models, ``can they be too
big?''~\cite{bender21}. For much of SE analytics, the honest answer is
yes. In one study, 52\% of ChatGPT answers to StackOverflow questions
were partly incorrect, and readers overlooked those errors 39\% of
the time~\cite{kabir24}. Scale the generation and you scale the
verification queue~\cite{baltes26}.

Some domains do not get a choice. The Green AI literature documents
the energy and carbon cost of large models~\cite{schwartz20,
strubell19}. Beyond that: much of the global south
lacks American-scale infrastructure; community colleges and K--12
classrooms teach on personal hardware; defense, healthcare, finance,
and law restrict what may leave the building; air-gapped and
on-device settings have no cloud at all. For all of these, easier AI
is the entry ticket.

SE itself has been teaching this lesson for decades. Modern software
is mostly assembled from pre-made parts, and more reuse means more
attack surface: left-pad, MOVEit (\$12 billion), and similar
supply-chain failures cost an estimated \$24 billion in 2025 alone.
Survivability argues the same way; SQLite ships its own tiny B-tree
rather than adopting Berkeley DB (which Oracle later acquired and
relicensed), and now runs on Mars. As SQLite's author puts it, ``if
you want to be free, that means doing things
yourself''~\cite{hipp21}. More code
is more to maintain and more to go wrong. Even our configuration
files are bloated: Xu et al.\ found that most users cannot cope with
most of the options we give them~\cite{xu15}.

There is also a supply problem. The big-data era ran on a silent
premise: more data exists. Villalobos et
al.\ project that the stock of human-generated public text runs out
within years~\cite{villalobos24}. Synthetic replacements trade away
privacy or fidelity~\cite{ling24}, and training recursively on
machine output collapses model quality~\cite{shumailov24}.

Finally, more was never always better. For agile effort estimation, a
support vector machine over TF-IDF features out-predicts deep learners
while running 100 times faster~\cite{tawosi23}. For text mining
StackOverflow, tuned simple learners ran 500 times faster than deep
learning, at no loss~\cite{majumder18}. Tree ensembles still beat deep
networks on typical tabular data~\cite{grinsztajn22}.
Often, better data or better tuning beats better (bigger)
learners~\cite{agrawal18, fu16, fu17}. Hooker calls this the hardware
lottery: ideas win because the available hardware favors them, not
because they are best~\cite{hooker21}. Big AI won a lottery. That is
not the same as winning the argument.

Summing up: the bitter lesson has preconditions (a clear objective,
abundant data, ever-cheaper compute). Chess and vision met all three;
much of SE meets none. Where the preconditions fail, acting as if the
lesson still applied is not rigor. It is habit.

\section{How: Active Learning with EZR}
\label{sec:how}
Less is a very old idea. Around
1300, William of Ockham warned against multiplying entities beyond
necessity. Pearson reduced data to a few principal components in
1901~\cite{pearson01}. Chang shrank nearest-neighbor classifiers to a
few prototypes in 1974~\cite{chang74}. Kohavi and John pruned feature
sets in 1997~\cite{kohavi97}. Active learning~\cite{settles09} and
model distillation continue the tradition. Of these, we focus on
active learning,
since labels are what SE data lacks. As Table~\ref{tab:xy} shows, for
$y=f(x)$ problems in SE, the independent features $x$ come cheap
while the dependent labels $y$ come dear: measuring runtimes, running
test suites, or eliciting expert judgment is slow and expensive,
while unlabeled examples are nearly free.

\begin{table}[t]
\caption{In SE, for $y=f(x)$, acquiring features $x$ is often far
cheaper than determining labels $y$.}
\label{tab:xy}
\centering
{\footnotesize
\begin{tabular}{p{.44\columnwidth}|p{.44\columnwidth}}
\hline
{\bf cheap feature acquisition ($x$)} & {\bf costly label determination ($y$)} \\
\hline
mine GitHub metadata (code size, dependencies) &
estimate market value or total development effort \\
count system classes &
audit actual maintenance effort \\
list design option permutations &
validate configurations via stakeholder consensus \\
enumerate software parameters &
execute full test suites for every unique build \\
identify learner hyper-parameters &
run exhaustive grid searches \\
generate fuzzing inputs &
execute tests, audit outputs \\
\hline
\end{tabular}}
\end{table}

Nor do the usual shortcuts fix this. Expert labeling is sluggish and
fatigue-prone, taking hours per case for security or technical-debt
judgments. Labels lifted from historical logs can mislead; in one
audit, most logged technical-debt entries were false
positives~\cite{yu22}. And asking an LLM to do the labeling does not
close the gap, since human-model agreement is no proxy for ground
truth~\cite{ahmed25}. Whatever the labeling method, budgets stay
small. So the research question becomes: how to spend a small budget
well?

The loop is short. Label a few random examples. Build a small model.
Use that model to pick the most informative unlabeled example: explore
where the model is most uncertain, exploit where it predicts the best
outcomes, and shift slowly from explore to exploit. Label it, rebuild,
repeat until the budget runs out. State-of-the-art optimizers such as
SMAC3 run the same loop over an ensemble of decision
trees~\cite{lindauer22}, and earlier SE work did so over pairs of
performance models~\cite{nair18}. But note the hidden cost in the
heavyweight versions: every new label triggers a rebuild of the whole
ensemble, so the machinery around the labels can cost far more than
the labels themselves. The EZR toolkit strips the loop to
a few hundred lines of dependency-free Python: a Naive Bayes-style
learner picks the labels, then a small regression tree summarizes
them~\cite{menzies26}. In that loop, labeling is the only expensive
step; everything else runs in milliseconds.

Testing such claims needs many tasks, not one case study. The MOOT
repository holds 127 multi-objective optimization datasets extracted
from recent SE papers~\cite{moot26}. The tasks cover software configuration, performance
tuning, product lines, project health, defect prediction, testing,
cost estimation, cross-domain generalization, and text mining. They
range from 3 to 1,044 independent attributes, one to eight goals, and
93 to 100,000 rows. Performance is scored as
$\mathit{win}=100\times(1-\mathit{regret})$, where regret is the
distance from the best result in the full dataset, normalized so that
100 means best-known and 0 means as bad as random guessing.

\begin{figure}[t]
\begin{scriptsize}
\begin{verbatim}
 win     n     Lbs-    Acc+    Mpg+
   8    43   2334.2    16.3    27.4
  41    26   2038.7    17.1    30.0  Volume <= 111
  71    10   1871.7    18.1    32.0  |  Volume <= 83
  85     4   1831.7    19.0    32.5  |  |  Model <= 73
  62     6   1898.3    17.4    31.7  |  |  Model > 73
  22    16   2143.0    16.5    28.7  |  Volume > 83
  36     7   2098.1    18.0    28.6  |  |  Model <= 73
  12     9   2177.9    15.3    28.9  |  |  Model > 73
  26     4   2210.0    15.7    30.0  |  |  |  Volume > 97
   0     5   2152.2    15.0    28.0  |  |  |  Volume <= 97
 -41    17   2786.1    15.1    23.5  Volume > 111
  -9     5   2279.8    14.6    28.0  |  Volume <= 116
 -54    12   2997.1    15.3    21.7  |  Volume > 116
 -44     8   2921.1    15.8    22.5  |  |  Model > 73
 -38     4   2563.0    14.0    25.0  |  |  |  Clndrs <= 4
 -50     4   3279.2    17.4    20.0  |  |  |  Clndrs > 4
 -75     4   3149.0    14.4    20.0  |  |  Model <= 73
\end{verbatim}
\end{scriptsize}
\caption{An EZR tree for a three-goal automotive task: minimize
weight (Lbs), maximize acceleration (Acc) and miles per gallon
(Mpg). Each line is one node: $\mathit{win}$ score, row count $n$,
mean goal values, and the branch test. Built from 43 labels; the
whole model fits in 17 lines and mentions three attributes.}
\label{fig:tree}
\end{figure}

This is best shown by example. In Figure~\ref{fig:tree}, rows
describe cars; the goals are to minimize weight while maximizing
acceleration and miles per gallon. From 43 labels, the root scores
$\mathit{win}=8$: barely better than guessing. Two tests down (engine
volume at most 83, model year at most '73) sit four cars scoring
$\mathit{win}=85$: light, quick, thrifty. The worst leaf
($\mathit{win}=-75$) holds heavy gas-guzzlers, so the tree also says,
just as plainly, what to avoid. The whole policy mentions three
attributes; a domain expert can dispute it line by line. To apply it,
sort new unlabeled examples down its branches and spend the labeling
budget only on those landing in the best leaves.

Does it work at scale? A study of 100,000 runs over MOOT proceeded as
follows~\cite{ganguly26}. Pick one of the 127 datasets. Hide the
labels; split off a holdout. Let the active learner spend at most 50
labels on the training side, then build a tree from what it learned.
Sort the holdout through that tree and label only the first five
suggestions, scoring the best of those five. Median result across all
runs: 85\% of the best known. The whole study ran in two hours on one
laptop.

How does that compare? In a tournament run at matched budgets,
NSGA-II~\cite{deb02} needed roughly 1,000 evaluations to reach what
these methods reach in 50~\cite{ganguly26b}, and EZR's per-task
runtimes sit in the milliseconds. That is not just faster science;
500 times less CPU is 500 times less energy, every run. Nor is this
the first such warning~\cite{chen19, menzies21}.

Are the resulting models small enough to trust? Across the 127
datasets, EZR's trees use four to nine attributes, out of as many as
1,044~\cite{rayegan26}. Systems described by a thousand variables
are routinely controlled by fewer than ten; effort spent modeling
the rest buys noise. Small is only useful if nothing important got
thrown away, so that work also ran a harder test: if EZR's tree uses
$N$ attributes, ask state-of-the-art attribute selectors for their
top $N$, build models from each selection, and compare. Trees built
from EZR's few attributes usually tied with, or came close to, the
best of those selectors~\cite{rayegan26}. A model that fits on half
a page can be checked by a domain expert before lunch. None of the
black-box alternatives offer that.

\section{Why Not More Easier AI?}
\label{sec:whynot}
If the evidence is this good, why is easier AI so rare? We offer four
answers.

First, weak empirical standards. LLM papers now flood SE
venues~\cite{hou24}, yet in a survey of 229 of them, only 5\% compared
against a simpler approach~\cite{menzies26}. That is an error, not a
style choice, since simpler methods can be better, faster, or
both~\cite{fu17, majumder18, tawosi23}. The fix is old: Cohen's 1995
empirical-AI text recommended baselining against straw-man
ablations~\cite{cohen95}. So, before shipping a clever method, build
the stupidest straw man that could possibly work. Surprisingly
often, the straw man wins, and the clever method should be thrown
away~\cite{menzies21}.

Second, incentives. Big sells. A data center has a ribbon the mayor
can cut. A 300-line script has no launch event, no special hardware to
sell (it is too fast to need any), and no valuation story. The
spectacle funds itself; the simplification does not.

Third, the simplicity paradox. Simplicity must be earned twice.
Earned in experience, since knowing what to cut takes years of
mistakes (``experience is the name every one gives to their
mistakes,'' as Wilde had it). And earned in compute, since certifying
that simple suffices means running both the simple and the complex
method at scale; the tournament behind \S\ref{sec:how} consumed some
20 months of CPU~\cite{ganguly26b}. Simple results are expensive to
prove and cheap-looking once proved. A bad bargain for a career, and
researchers know it.

Fourth, our brains fight us. In studies covering hundreds of
subjects, Adams et al.\ found that people systematically reach for
additive changes over subtractive ones, by roughly four to
one~\cite{adams21}. We suspect the simpler methods reported here were
not found earlier for the simplest of reasons: nobody was looking
down.

Call the sum of these the inverse bitter lesson. Sutton warned that
researchers cling to clever hand-crafted methods when scale would do
better~\cite{sutton19}. The companion failure is just as real:
reaching for scale reflexively, an LLM and an embedding API for a
problem that thirty lines of Naive Bayes over 1,700 rows would solve.
The first failure wastes a project. The second, repeated across 95\%
of a literature~\cite{menzies26}, wastes a field.

\section{Future Plans}
\label{sec:future}
The emerging result: 50 labels often suffice. The new idea: treat
that as a research program, the data-light challenge, asking for
which SE tasks a handful of labels suffices and when it
fails~\cite{ganguly26}. A full-length result needs four things.

Characterization. Not all tasks simplify; we need to predict which.
The plan: regress label-budget requirements against measurable task
features (dimensionality, goal count, noise, intrinsic dimension)
across an expanded MOOT~\cite{moot26} that adds failure cases, since
a corpus that simple methods always win teaches nothing.

Tournaments. The comparison set must include what people actually
use: state-of-the-art optimizers~\cite{lindauer22, deb02} and
LLM-based agents, run head-to-head at matched labeling
budgets~\cite{ganguly26b}, with energy reported alongside accuracy.

Hybrids. Easier AI need not be anti-LLM. Early results suggest LLMs
can warm-start active learning with a few good initial
guesses~\cite{senthilkumar26}, and that running classical optimizers
before LLMs beats the reverse order~\cite{srinivasan26}. The open
question is where the crossover sits: how few labels before the prior
knowledge inside an LLM stops paying rent?

Standards and teaching. We will draft a one-line addition to SE
review checklists (``compared to what simpler?'') and an
easier-AI-first curriculum, where students meet the 300-line version
before the trillion-parameter one.

What could kill this idea? Three things, all testable.
MOOT could be biased toward tasks where simple methods win
(fix: grow the corpus adversarially). Safety-critical work may
reject 85\% of best (the characterization study should mark that
boundary). And the evidence here is tabular; whether it extends to
code, traces, and text is what the hybrid studies will answer. Either
way, the community gets what it lacks: an evidence-based boundary
between easy and hard.

\section{Conclusion}
Lehman told us that complexity grows unless work is done to reduce
it~\cite{lehman80}. This paper argues we must (economics, energy, trust, access) and
that we can (50 labels, four to nine attributes, 85\% of best, one
laptop). The tools and data are public~\cite{menzies26, moot26}.
Sutton was right where feedback is free. The sweeter lesson covers
the rest: where labels are dear, a few dozen well-chosen ones go
most of the way. Baseline the simple. At review time, ask
``compared to what simpler?'' Teach the easy version first. Scale
only when the simple fails.

%\section*{Acknowledgements}
% Camera-ready only (double-anonymous). Students, per the supervision-tree
% slide of the talk: Kishan Kumar Ganguly, Amirali Rayegan, Srinath
% Srinivasan, Lohith Senthilkumar, Amritanshu Agrawal, Vivek Nair,
% Jianfeng Chen, Andre Lustosa.

\balance
\bibliographystyle{IEEEtran}
\bibliography{nier}

\end{document}
